\chapter{补充材料}

\section{Balance Loss 计算流程}

表~\ref{tab:balance_loss_flow}给出了 Balance Loss 框架中 Inter-CBL 和 Intra-CBL 的计算流程。

\begin{table}[htbp]
\centering
\caption{Balance Loss 计算流程}
\label{tab:balance_loss_flow}
\begin{tabular}{p{0.88\textwidth}}
\hline
\textbf{流程}：Balance Loss 计算流程 \\
\hline
\textbf{输入}：预测概率图 $\mathbf{p} \in [0,1]^{H \times W}$，真实标签 $\mathbf{y} \in \{0,1\}^{H \times W}$，\\
\quad\quad\quad 超参数 $\alpha, \beta, \gamma, \tau, w_h, \text{neg\_ratio}$ \\
\textbf{输出}：总损失 $L_{balance}$ \\
\hline
\textbf{Step 1: Inter-CBL（类别间平衡损失）} \\
1. 分离前景像素集合 $F = \{i : y_i = 1\}$ 和背景像素集合 $B = \{i : y_i = 0\}$ \\
2. 计算所有背景像素的前景预测概率 $\{p_j\}_{j \in B}$ \\
3. 按前景概率降序排列背景像素，取前 $|F| \times \text{neg\_ratio}$ 个作为困难背景集合 $H_B$ \\
4. 计算 $L_{inter} = \frac{1}{2}\left(\frac{1}{|F|}\sum_{i \in F} \ell_i + \frac{1}{|H_B|}\sum_{j \in H_B} \ell_j\right)$ \\
\\
\textbf{Step 2: Intra-CBL（类别内难度加权损失）} \\
5. 计算每个像素的置信度误差 $e_i = |p_i - y_i|$ \\
6. 按阈值 $\tau$ 划分：困难像素 $H = \{i : e_i > \tau\}$，易分像素 $E = \{i : e_i \leq \tau\}$ \\
7. 计算 $L_{intra} = 1.0 \cdot \frac{1}{|E|}\sum_{i \in E} \ell_i + w_h \cdot \frac{1}{|H|}\sum_{j \in H} \ell_j$ \\
\\
\textbf{Step 3: 损失组合} \\
8. 计算 Dice 损失 $L_{dice} = 1 - \frac{2|\mathbf{y} \cap \hat{\mathbf{y}}|}{|\mathbf{y}| + |\hat{\mathbf{y}}|}$ \\
9. \textbf{return} $L_{balance} = \alpha \cdot L_{inter} + \beta \cdot L_{intra} + \gamma \cdot L_{dice}$ \\
\hline
\end{tabular}
\end{table}

\section{两阶段训练策略}

\begin{itemize}
  \item \textbf{Stage 1}（Epoch 0 至 $T_1=50$）：仅启用 $L_{intra} + L_{dice}$，此阶段 Inter-CBL 不参与计算。目的是在模型预测近似随机的初期避免困难背景挖掘引入噪声梯度。
  \item \textbf{Stage 2}（Epoch $T_1$ 至终止）：引入完整的 $\alpha L_{inter} + \beta L_{intra} + \gamma L_{dice}$。此时模型已具备基础判别能力，Inter-CBL 的困难背景挖掘可以产生有效的梯度信号。
\end{itemize}

$T_1=50$ 的选择基于初步实验观察：在前 50 个 epoch 内，模型预测逐渐从随机状态收敛至基础判别能力，此后引入 Inter-CBL 可避免早期噪声梯度的干扰。

\section{局部适配器核心代码}

以下为局部适配器模块的 PyTorch 实现代码：

\begin{lstlisting}
import torch
import torch.nn as nn

class MSLAdapter(nn.Module):
    """Multi-Scale Local Adapter: dual-path depthwise separable
    convolution for multi-scale local feature enhancement."""

    def __init__(self, embed_dim=256):
        super().__init__()
        # Path 1: standard 3x3 depthwise conv (local)
        self.conv1 = nn.Conv2d(
            embed_dim, embed_dim, kernel_size=3,
            padding=1, groups=embed_dim, bias=True
        )
        # Path 2: dilated 3x3 depthwise conv
        #          (larger local receptive field)
        self.conv2 = nn.Conv2d(
            embed_dim, embed_dim, kernel_size=3,
            padding=2, dilation=2, groups=embed_dim, bias=True
        )
        # Pointwise conv for channel fusion (512 -> 256)
        self.pointwise = nn.Conv2d(
            embed_dim * 2, embed_dim, kernel_size=1, bias=True
        )
        self.act = nn.GELU()

    def forward(self, x):
        # x: (B, 256, 64, 64) from Image Encoder
        f1 = self.act(self.conv1(x))   # local features
        f2 = self.act(self.conv2(x))   # larger-scale local features
        f_cat = torch.cat([f1, f2], dim=1)  # (B, 512, 64, 64)
        f_out = self.pointwise(f_cat)       # (B, 256, 64, 64)
        return x + f_out  # residual connection
\end{lstlisting}

\textbf{参数量统计}：conv1 = 2,560，conv2 = 2,560，pointwise = 131,328，总计 \textbf{136,448} 个参数（占 MedSAM 93.7M 总参数的 0.15\%）。

\section{超参数配置汇总}

表~\ref{tab:hyperparams}汇总了本文各实验方案的关键超参数配置。

\begin{table}[htbp]
  \centering
  \caption{各实验方案超参数配置}
  \label{tab:hyperparams}
  \footnotesize
  \setlength{\tabcolsep}{2.5pt}
  \begin{tabular}{lccccccc}
    \hline
    参数 & Baseline & Intra-$\tau$0.9 & BL & BL+CA & BL+LoRA-r4 & BL+LoRA-r8 & BL+MSL \\
    \hline
    \multicolumn{8}{l}{\textit{共享训练配置$^*$}} \\
    优化器 & \multicolumn{7}{c}{AdamW（lr=$10^{-4}$, weight\_decay=0.01）} \\
    训练轮数 / Batch size & \multicolumn{7}{c}{200 epochs / 8} \\
    学习率调度 & \multicolumn{7}{c}{余弦退火} \\
    数据增强 & \multicolumn{7}{c}{H-Flip($p$=0.5) + V-Flip($p$=0.5) + Rot90} \\
    提示扰动 & \multicolumn{7}{c}{GT box 四边各 $[0, 20]$ px 随机外扩（仅训练阶段）} \\
    随机种子 & \multicolumn{7}{c}{2025} \\
    评估 checkpoint & \multicolumn{7}{c}{best（训练损失最低的 epoch）} \\
    \hline
    $\alpha$ (Inter-CBL) & -- & -- & 0.5 & 0.5 & 0.5 & 0.5 & 0.5 \\
    $\beta$ (Intra-CBL) & -- & 1.0 & 1.0 & 1.0 & 1.0 & 1.0 & 1.0 \\
    $\gamma$ (Dice) & -- & 1.0 & 1.0 & 1.0 & 1.0 & 1.0 & 1.0 \\
    $\tau$ (threshold) & -- & 0.9 & 0.9 & 0.9 & 0.9 & 0.9 & 0.9 \\
    $w_h$ (hard weight) & -- & 2.0 & 2.0 & 2.0 & 2.0 & 2.0 & 2.0 \\
    neg\_ratio & -- & -- & 3.0 & 3.0 & 3.0 & 3.0 & 3.0 \\
    stage1\_epochs & -- & -- & 50 & 50 & 50 & 50 & 50 \\
    \hline
    LoRA rank & -- & -- & -- & -- & 4 & 8 & -- \\
    LoRA $\alpha_{lora}$ & -- & -- & -- & -- & 1.0 & 1.0 & -- \\
    主干冻结$^\dag$ & 否 & 否 & 否 & 否 & 是 & 是 & 否 \\
    Cross-Attention & 否 & 否 & 否 & 是 & 否 & 否 & 否 \\
    局部适配器 & 否 & 否 & 否 & 否 & 否 & 否 & 是 \\
    \hline
  \end{tabular}
  \vspace{2pt}
  \parbox{\textwidth}{\footnotesize $^*$所有方案共享上述训练配置，仅在损失参数与微调方式上存在差异。\\
  $^\dag$“主干冻结”指 Image Encoder 参数冻结；Prompt Encoder 在所有方案中均保持冻结。BL+LoRA-r16（rank=16）配置与 r4/r8 一致，因表宽限制未单独列出。BL+CA 的额外参数来自 Cross-Attention 模块（$\sim$263K），推理时需额外编码参考样本。}
\end{table}

\section{实验环境}

\begin{itemize}
  \item \textbf{GPU}：4$\times$ NVIDIA V100（32GB 显存）
  \item \textbf{典型训练时间}：BL+MSL 约 52 小时/200 epochs（详见第5章表~\ref{tab:training_overhead}）
  \item \textbf{深度学习框架}：PyTorch 2.0
  \item \textbf{操作系统}：Ubuntu 20.04 LTS
  \item \textbf{Python 版本}：3.10
  \item \textbf{CUDA 版本}：11.8
  \item \textbf{主要依赖库}：numpy 1.24, scipy 1.10, monai 1.2, opencv-python 4.7
  \item \textbf{基座模型}：MedSAM（基于 ViT-Base，93.7M 参数）
  \item \textbf{预训练权重}：MedSAM 官方发布权重\cite{ma2024medsam}
\end{itemize}

\section{器官列表}

本文实验涉及的 13 个腹部器官如下：

\begin{itemize}
  \item \textbf{实质器官}（6 个）：肝脏（Liver）、脾脏（Spleen）、左肾（Left Kidney）、右肾（Right Kidney）、胰腺（Pancreas）、胆囊（Gallbladder）
  \item \textbf{空腔器官}（3 个）：食管（Esophagus）、胃（Stomach）、十二指肠（Duodenum）
  \item \textbf{血管}（2 个）：主动脉（Aorta）、下腔静脉（Inferior Vena Cava, IVC）
  \item \textbf{内分泌器官}（2 个）：右侧肾上腺（Right Adrenal Gland, RAG）、左侧肾上腺（Left Adrenal Gland, LAG）
\end{itemize}

\textbf{说明}：本文使用的 AMOS22 数据集原始标注含 15 类器官，已排除 Bladder（膀胱）与 Prostate/Uterus（前列腺/子宫）后保留 13 类，与 FLARE22 的 13 类标签体系一致。

\section{代码与数据可用性说明}

本文所涉及的 Balance Loss 实现代码、MSL-Adapter 模块定义、训练与评估脚本均基于 MedSAM\cite{ma2024medsam}开源代码库进行开发。数据集方面，本文使用的 AMOS22 CT 子集可通过官方渠道申请获取。
